% 神经网络章节的可复用图形语法；颜色在 layout.tex 集中定义。
\tikzset{
  nn neuron/.style={circle,minimum size=8mm,inner sep=0pt,
    line width=1.1pt,text=NNInk,draw=NNHidden,fill=NNHidden!18},
  nn input/.style={nn neuron,draw=NNInput,fill=NNInput!18},
  nn hidden/.style={nn neuron,draw=NNHidden,fill=NNHidden!18},
  nn output/.style={nn neuron,draw=NNOutput,fill=NNOutput!18},
  nn link/.style={draw=NNLine,line width=.85pt},
  nn flow/.style={-{Stealth[length=2mm,width=1.4mm]},
    draw=NNInput,line width=1.2pt},
  nn gradient/.style={-{Stealth[length=2mm,width=1.4mm]},
    draw=NNGradient,line width=1.2pt,dashed},
  nn heading/.style={font=\figurefont\bfseries\footnotesize,text=NNInk},
  nn annotation/.style={font=\figurefont\footnotesize,text=BookMuted},
  nn transform/.style={rounded corners=1.5mm,draw=NNOutput,
    fill=NNOutput!18,line width=1.1pt,minimum height=9mm,
    inner sep=3pt,align=center,text=NNInk}
}

% 操作示例的数值网格：左上 x/y、行数、列数、格宽、边框色、行/列/值。
% 空间索引从零开始，色彩和字号沿用神经网络图形语法。
\newcommand{\cnnGrid}[7]{%
  \draw[draw=#6,line width=.95pt] (#1,#2)
    rectangle ({#1+#4*#5},{#2-#3*#5});%
  \ifnum#4>1\relax
  \pgfmathtruncatemacro{\cnnGridLast}{#4-1}%
  \foreach \cnnGridIndex in {1,...,\cnnGridLast}{%
    \draw[nn link,line width=.5pt] ({#1+\cnnGridIndex*#5},#2)--
      ({#1+\cnnGridIndex*#5},{#2-#3*#5});%
  }%
  \fi
  \ifnum#3>1\relax
  \pgfmathtruncatemacro{\cnnGridLast}{#3-1}%
  \foreach \cnnGridIndex in {1,...,\cnnGridLast}{%
    \draw[nn link,line width=.5pt] (#1,{#2-\cnnGridIndex*#5})--
      ({#1+#4*#5},{#2-\cnnGridIndex*#5});%
  }%
  \fi
  \foreach \cnnGridRow/\cnnGridCol/\cnnGridValue in {#7}{%
    \node at ({#1+(\cnnGridCol+.5)*#5},{#2-(\cnnGridRow+.5)*#5})
      {$\cnnGridValue$};%
  }%
}

% CNN 的长方体表示空间与通道；局部块、角点连线及平面 MLP 共用颜色。
\pgfdeclarelayer{cnn background}
\tikzset{
  cnn plane/.style={draw=NNHidden,line width=1.1pt},
  cnn flow/.style={nn flow},
  cnn skip/.style={nn flow,draw=NNHidden,line width=1.1pt,rounded corners=2mm},
  cnn mapping/.style={draw=NNHidden!48,line width=.85pt,
    dash pattern=on 2.2pt off 1.5pt},
  cnn scale/.style={draw=NNLine,line width=.85pt},
  cnn title/.style={font=\figurefont\bfseries\footnotesize,text=NNInk,align=center},
  cnn size/.style={font=\figurefont\footnotesize,text=NNInk,align=center},
  cnn operation/.style={font=\figurefont\footnotesize,text=BookMuted,align=center},
  cnn neuron/.style={nn hidden,minimum size=3.5mm},
  cnn merge/.style={circle,draw=NNInk,fill=NNCanvas,
    minimum size=5.5mm,inner sep=0pt,line width=1.1pt}
}
% 单行架构的统一标题线、尺寸线；真正并行的分支可保留局部标题。
\newcommand{\cnnArchitecture}[2]{%
  \def\cnnTitleY{#1}\def\cnnSizeY{#2}%
  \pgfsetlayers{cnn background,main}%
}
\newcommand{\cnnLayerTitle}[3]{%
  \ifdefined\cnnTitleY
    \node[cnn title,anchor=north] at ({#1},\cnnTitleY) {#3};%
  \else
    \node[cnn title,anchor=south] at ({#1},{#2}) {#3};%
  \fi
}
\newcommand{\cnnLayerSize}[3]{%
  \ifdefined\cnnSizeY
    \node[cnn size,anchor=north] at ({#1},\cnnSizeY) {#3};%
  \else
    \node[cnn size,anchor=north] at ({#1},{#2}) {#3};%
  \fi
}
% 名称、外框左下 x/y、正面宽/高、示意厚度、颜色、标题、实际形状。
\newcommand{\cnnFeature}[9]{%
  \pgfmathsetmacro{\cnnDepth}{.7*#6}%
  \colorlet{cnn-#1-color}{#7}%
  \draw[cnn plane,draw=#7,fill=#7!12]
    (#2,{#3+#5+\cnnDepth})--({#2+#4},{#3+#5+\cnnDepth})--
    ({#2+#4+\cnnDepth},{#3+#5})--({#2+\cnnDepth},{#3+#5})--cycle;%
  \draw[cnn plane,draw=#7,fill=#7!30]
    (#2,{#3+\cnnDepth})--({#2+\cnnDepth},#3)--
    ({#2+\cnnDepth},{#3+#5})--(#2,{#3+#5+\cnnDepth})--cycle;%
  \draw[cnn plane,draw=#7,fill=#7!20]
    ({#2+\cnnDepth},#3) rectangle ({#2+#4+\cnnDepth},{#3+#5});%
  \coordinate (#1-west) at (#2,{#3+#5/2+\cnnDepth/2});%
  \coordinate (#1-east) at ({#2+#4+\cnnDepth},{#3+#5/2});%
  \coordinate (#1-north) at ({#2+(#4+\cnnDepth)/2},{#3+#5+\cnnDepth});%
  \coordinate (#1-south) at ({#2+(#4+\cnnDepth)/2},#3);%
  % 相邻张量连接右侧面与左侧面；分别保留前后、上下四个角点。
  \foreach \cnnLevel/\cnnY in {bottom/0,top/1}{%
    \coordinate (#1-left-front-\cnnLevel) at ({#2+\cnnDepth},{#3+\cnnY*#5});%
    \coordinate (#1-left-back-\cnnLevel) at (#2,{#3+\cnnY*#5+\cnnDepth});%
    \coordinate (#1-right-front-\cnnLevel) at ({#2+#4+\cnnDepth},{#3+\cnnY*#5});%
    \coordinate (#1-right-back-\cnnLevel) at ({#2+#4},{#3+\cnnY*#5+\cnnDepth});%
  }%
  \foreach \cnnCorner/\cnnX/\cnnY in {sw/0/0,se/1/0,nw/0/1,ne/1/1}{%
    \coordinate (#1-face-\cnnCorner) at ({#2+\cnnDepth+\cnnX*#4},{#3+\cnnY*#5});%
    \coordinate (#1-patch-\cnnCorner) at ({#2+\cnnDepth+(.18+.37*\cnnX)*#4},{#3+(.22+.31*\cnnY)*#5});%
    \coordinate (#1-patch-back-\cnnCorner) at ({#2+(.18+.37*\cnnX)*#4},{#3+(.22+.31*\cnnY)*#5+\cnnDepth});%
    \coordinate (#1-cell-\cnnCorner) at ({#2+\cnnDepth+(.2+.14*\cnnX)*#4},{#3+(.32+.14*\cnnY)*#5});%
    \coordinate (#1-cell-back-\cnnCorner) at ({#2+(.2+.14*\cnnX)*#4},{#3+(.32+.14*\cnnY)*#5+\cnnDepth});%
  }%
  \foreach \cnnUnit/\cnnRatio in {1/.2,2/.5,3/.8}{%
    \coordinate (#1-vout-\cnnUnit) at ({#2+#4+\cnnDepth},{#3+\cnnRatio*#5});%
    \coordinate (#1-vin-\cnnUnit) at ({#2+\cnnDepth},{#3+\cnnRatio*#5});%
  }%
  \cnnLayerTitle{#2+(#4+\cnnDepth)/2}{#3+#5+\cnnDepth+2.5}{#8}%
  \cnnLayerSize{#2+(#4+\cnnDepth)/2}{#3-2.5}{#9}%
}
\newcommand{\cnnScaleLinks}[2]{%
  \begin{pgfonlayer}{cnn background}%
  \foreach \cnnSide in {front,back}{\foreach \cnnLevel in {bottom,top}{%
    \draw[cnn scale] (#1-right-\cnnSide-\cnnLevel)--(#2-left-\cnnSide-\cnnLevel);%
  }%
  }%
  \end{pgfonlayer}%
}
% 小块沿通道厚度贯穿，以所属张量的深色突出。
\newcommand{\cnnLocalBlock}[2][NNHidden]{%
  \draw[draw=#1!80!NNInk,fill=#1!38,line width=1.1pt]
    (#2-back-nw)--(#2-back-ne)--(#2-ne)--(#2-nw)--cycle;%
  \draw[draw=#1!80!NNInk,fill=#1!65,line width=1.1pt]
    (#2-back-sw)--(#2-sw)--(#2-nw)--(#2-back-nw)--cycle;%
  \draw[draw=#1!80!NNInk,fill=#1!48,line width=1.1pt]
    (#2-sw)--(#2-se)--(#2-ne)--(#2-nw)--cycle;%
}
\newcommand{\cnnMap}[2]{%
  % 局部块同样只连接彼此相向的侧面：右侧面的四角对应左侧面的四角。
  \foreach \cnnFrom/\cnnTo in {se/sw,ne/nw,back-se/back-sw,back-ne/back-nw}{%
    \draw[cnn mapping] (#1-patch-\cnnFrom)--(#2-cell-\cnnTo);%
  }%
  \cnnLocalBlock[cnn-#1-color]{#1-patch}%
  \cnnLocalBlock[cnn-#2-color]{#2-cell}%
}
\newcommand{\cnnWholeLinks}[2]{%
  \foreach \cnnFrom in {1,2,3}{\foreach \cnnTo in {1,2,3}{%
    \draw[nn link] (#1-vout-\cnnFrom)--(#2-vin-\cnnTo);%
  }}%
}
% 平面全连接层：可选颜色、名称、中心 x/y、层名、单元数。
\newcommand{\cnnFC}[6][NNHidden]{%
  \draw[draw=#1,fill=#1!8,line width=1.1pt]
    ({#3-3.4},{#4-13}) rectangle ({#3+3.4},{#4+13});%
  \foreach \cnnUnit/\cnnY in {1/-8,2/0,3/8}{%
    \node[cnn neuron,draw=#1,fill=#1!18] (#2-\cnnUnit) at (#3,{#4+\cnnY}) {}; %
    \coordinate (#2-vin-\cnnUnit) at ({#3-1.75},{#4+\cnnY});%
    \coordinate (#2-vout-\cnnUnit) at ({#3+1.75},{#4+\cnnY});%
  }%
  \coordinate (#2-west) at ({#3-3.4},#4);%
  \coordinate (#2-east) at ({#3+3.4},#4);%
  \cnnLayerTitle{#3}{#4+15.5}{#5}%
  \cnnLayerSize{#3}{#4-15.5}{#6}%
}
\newcommand{\cnnDenseLinks}[2]{\cnnWholeLinks{#1}{#2}}

% 平面计算图采用浅色圆角模块、简短算子标签和连贯的圆角跨层路径。
\tikzset{
  nn operator/.style={draw=NNHidden,fill=NNHidden!18,rounded corners=1.5mm,
    line width=1.1pt,minimum height=10mm,inner sep=0pt,align=center},
  nn operator path/.style={nn flow,draw=NNHidden,rounded corners=2.5mm},
  nn pointwise/.style={circle,draw=NNHidden,fill=NNCanvas,line width=1.1pt,
    minimum size=6mm,inner sep=0pt}
}
% 名称、中心 x/y、核尺寸。
\newcommand{\nnConvOp}[4]{%
  \node[nn operator,fill=NNHidden!18,minimum width=10mm,minimum height=13mm]
    (#1) at (#2,#3) {Conv\\$#4$};%
}
\newcommand{\nnReluOp}[3]{%
  \node[nn operator,fill=NNHidden!18,minimum width=10mm] (#1) at (#2,#3) {ReLU};%
}
\newcommand{\nnSigmoidOp}[3]{%
  \node[nn operator,fill=NNHidden!18,minimum width=17mm] (#1) at (#2,#3) {Sigmoid};%
}
\newcommand{\nnBNOp}[3]{%
  \node[nn operator,fill=NNHidden!18,minimum width=8mm] (#1) at (#2,#3) {BN};%
}
% 可选圆直径、名称、中心 x/y；青色圆形 + 表示通道拼接，图注说明其语义。
\newcommand{\nnConcatOp}[4][6mm]{%
  \node[nn pointwise,draw=NNOutput,minimum size=#1]
    (#2) at (#3,#4) {$+$};%
}
\newcommand{\nnTensorState}[5]{%
  \node[draw=#4,fill=#4!8,line width=1.1pt,minimum width=8mm,
    minimum height=20mm,inner sep=0pt] (#1) at (#2,#3) {};%
  \foreach \nnOffset in {-6,0,6}{%
    \node[circle,draw=#4,fill=#4!20,line width=1.1pt,
      minimum size=3mm,inner sep=0pt] at (#2,{#3+\nnOffset}) {};%
  }%
  \node[nn heading,anchor=south] at (#2,{#3+12}) {#5};%
}


% 用于第18章旧图形组件的兼容样式；语义颜色与后续章节相同。
\tikzset{
  nn legacy node/.style={nn operator,inner sep=5pt},
  nn legacy arrow/.style={nn flow,draw=NNWire},
  nn frame/.style={draw=NNLine,fill=NNPanel,line width=.85pt,rounded corners=1.5mm}
}
